\section{Pricing the floor before releasing the margin}
\label{sec:insurance-priced-guarantee}

The settlement schedule identifies a liability that increases when the
underlying assets fall. To measure that liability before settlement, we need
a financial model. A small, explicitly specified market makes the distinction
between a price and a forecast visible. The retained standard permits a
replicating portfolio to measure cash flows that it matches in amount, timing
and uncertainty; other techniques must remain consistent with observable
market prices \citep[paragraphs B44--B48]{aasb17}.

Suppose a non-dividend-paying traded asset costs $S_0=100$. Its next observed
value is either $S_u=120$ or $S_d=80$, after which it remains at that value
until payment. Cash earns zero interest, so its accumulation factor is $R=1$.
There are exactly two financial states, unrestricted borrowing and shorting,
no transaction costs and no counterparty failure. These are assumptions of
the illustration. The financial observation precedes all supplied coverage
units; the subsequent service interval contains the mortality risk. An actual
period with service already provided would require its own dated allocation.

For payoffs $X_u,X_d$, let $\delta$ denote the units of the traded asset and
$b$ the amount deposited in cash. Matching the two payments requires
\[
 \delta S_u+bR=X_u,\qquad \delta S_d+bR=X_d.
\]
Subtracting the second equation from the first gives
$\delta=(X_u-X_d)/(S_u-S_d)$. Substitution then determines $b$ and the
portfolio's current cost:
\[
 b=\frac{S_uX_d-S_dX_u}{(S_u-S_d)R},\qquad
 V_0=\delta S_0+b
     =\frac{qX_u+(1-q)X_d}{R},\qquad
 q=\frac{RS_0-S_d}{S_u-S_d}.
\]
Thus the weight $q$ follows from today's asset and cash prices. Positive
state prices require $S_d<RS_0<S_u$, with $R>0$. Here $q=1/2$.
Additional financial states or restrictions that prevent the matching trades
would invalidate this particular valuation method.

The policyholder's investment payment is $\max(0.9S,90)$. Its top-up above
the ordinary share is $(90-0.9S)^+$, where $x^+=\max(x,0)$.
That top-up pays zero in the up state and 18 in the down state.
The matching portfolio has $\delta=-0.45$ and $b=54$:
$-0.45(120)+54=0$ and $-0.45(80)+54=18$.
Its cost is $-0.45(100)+54=9$. Adding the ordinary share's value of 90
gives a financial-benefit value of 99. This derives the option value; no
option price has been supplied as an unexplained input.

For a separate physical forecast, suppose the up state has probability
$3/4$. The expected top-up is then $18/4=4.5$, although its market-consistent
price remains 9. Suppose also that a covered death triggers an additional
payment of 20, with physical probability $1/10$, independently of the market
state. Its expected cost is 2. Nonfinancial risk adjustment of 1 is a
separately supplied measurement assumption, with its entire release in the
service result. The financial state prices, physical joint law and compensation
for nonfinancial risk therefore occupy different records
\citep[paragraphs 33, 37, 81 and B42--B49]{aasb17}.

\subsection{The same measurement through service and payment}

Use the retained historical Australian provisions for a conditional
calendar-2024/25 example. Assume the contract has significant insurance risk,
an enforceable link to the identified pool and the required direct-participation
characteristics. The component rights terminate together; neither can be
retained separately after the other matures. This supplies the specified
interrelationship under B32(b), together with an assessment that no embedded
derivative must be separated. These admission findings need actual terms
before real use; neither $q$ nor the 90\% share proves them
\citep[paragraphs 10--13, B17--B23, B31--B32 and B101--B108]{aasb17}.

The investment payment is due whether the insured lives or dies. There are
no surrender, lapse, acquisition or tax cash flows in this example, and all
figures are pretax. The insurer holds one unit of the underlying asset at
fair value through profit or loss and makes no risk-mitigation election.
The portfolio used to calculate the guarantee price is not its actual hedge.
Insurance finance also enters profit or loss under paragraph 89(a).
The investment component remains within the insurance measurement but is
excluded from service revenue and expense under paragraph 85
\citep[paragraphs 85, 89 and B115--B118]{aasb17}.

Start with owner capital and cash of 50, and receive premium 110.
Before receipt, fulfilment cash flows are $99+2+1-110=-8$, giving CSM 8.
After receipt, the remaining-coverage liability is 110: financial benefits 99,
mortality estimate 2, risk adjustment 1 and CSM 8. Buying the underlying
for 100 leaves cash 60 and investment 100 against liability 110 and equity 50.

At the financial observation, the future-service fee change is the change
in the insurer's 10\% share less the change in the top-up value
\citep[paragraph 45 and B111--B114]{aasb17}. In the up state this is
$2-(0-9)=11$, increasing CSM from 8 to 19. Fulfilment cash flows become
$108+2+1=111$ and the liability becomes 130. Investment gain 20 and
insurance finance expense 20 offset.

In the down state the fee change is $-2-(18-9)=-11$. The existing CSM
absorbs 8; the excess 3 becomes a service loss. Fulfilment cash flows and the
liability are $90+2+1=93$, with CSM zero and a loss component of 3
\emph{within} that liability. Investment loss 20 and insurance finance income
20 offset, while the service loss reduces equity to 47. The loss component
does not create a second liability, and the binding floor does not reopen
the inception VFA classification
\citep[paragraphs 45, 48--52, 87(c) and B102]{aasb17}.

Completing the service interval releases the expected mortality cost, risk
adjustment and remaining CSM. Up-state revenue is $2+1+19=22$.
In the down state the service-related release of 3 reverses the loss component,
leaving revenue zero. Actual mortality expense is either zero or 20, first
recognised through the incurred-claim liability and then paid. Selling the
underlying and paying the investment component of 108 or 90 completes all
balances.

\begin{center}
\begin{tabular}{@{}lrrrr@{}}
\toprule
Joint outcome & Physical weight & Profit & Final cash/equity & OCI\\
\midrule
Up, survival   & $27/40$ & 22  & 72 & 0\\
Up, death      & $3/40$  & 2   & 52 & 0\\
Down, survival & $9/40$  & 0   & 50 & 0\\
Down, death    & $1/40$  & $-20$ & 30 & 0\\
\bottomrule
\end{tabular}
\end{center}

Each row is a complete conditional history. Operating cash is premium 110
less the investment benefit and mortality payment; the investment purchase
and sale give investing cash of $+20$ or $-20$. This classification assumes
nontrading investments under the retained cash-flow provisions; fair-value
measurement alone does not determine trading purpose
\citep[paragraphs 14--16]{aasb107dated}. The physical expected profit is 14.5
under the supplied law, but the four rows preserve the separate downside
exposures. An unknown mortality law, risk adjustment, dependence assumption
or admission finding leaves a named unresolved measurement rather than a
default probability or zero cost.

The exact state prices, journals, component changes and adverse alternatives
are retained with the acquisition/treaty cases in
\path{data/accounting-master/}\allowbreak\path{finite-method-decisions-20261015/insurance/}.
The companion method, \emph{Pricing the guarantee and retaining separate
insurance groups}, binds valuation, probability and risk-adjustment records to the same group,
date and evidence view, while keeping their meanings distinct. General market
models and actuarial calibration require additional methods and validation.

